\documentclass[12pt]{article}
\usepackage{amsmath, amssymb, amsthm}
\usepackage{geometry}
\usepackage{enumitem}
\geometry{margin=1in}

\title{Analysis of \textbf{Claude's Cycles}: Validity of the Construction}
\author{}
\date{}

\begin{document}

\maketitle

\section{Introduction}
This document analyzes the algorithm proposed in Don Knuth's \textit{Claude's Cycles} for decomposing a specific directed graph into three Hamiltonian cycles. The graph consists of $m^3$ vertices $(i, j, k)$ with $0 \le i, j, k < m$. The edges allow incrementing any coordinate modulo $m$.

The paper claims that Claude Opus 4.6 discovered a general solution for all odd $m > 1$. However, the "simplification" of this solution provided in the C code snippet on page 2 contains a logical flaw that prevents it from generating valid Hamiltonian cycles for $c=1$ and $c=2$.

We will prove:
\begin{enumerate}
    \item The logic for $c=0$ generates a valid Hamiltonian cycle.
    \item The logic for $c=1$ and $c=2$ is mathematically incapable of visiting all $m^3$ vertices (it gets stuck in short loops).
    \item This does not mean the problem is unsolvable; rather, it means the \textbf{simplification} on page 2 is incorrect. The true solution (the Python program) must use a different, more complex logic.
\end{enumerate}

\section{The Algorithm Structure}
Let $m$ be an odd integer. The algorithm maintains a state $(i, j, k)$ and a cycle index $c \in \{0, 1, 2\}$.
At each step, it computes $s = (i + j + k) \pmod m$ and selects a permutation string $d \in \{\text{"012"}, \text{"120"}, \text{"210"}\}$ based on rules involving $s$, $i$, and $j$.
The next move for cycle $c$ is determined by the character at index $c$ of $d$:
\begin{itemize}
    \item $d[c] = \text{'0'} \implies i \leftarrow (i+1) \pmod m$
    \item $d[c] = \text{'1'} \implies j \leftarrow (j+1) \pmod m$
    \item $d[c] = \text{'2'} \implies k \leftarrow (k+1) \pmod m$
\end{itemize}

\section{Proof of Validity for $c=0$}
We show that the sequence generated for $c=0$ visits every vertex exactly once.

The rules for $c=0$ are:
\begin{itemize}
    \item If $s=0$:
    \begin{itemize}
        \item If $j=m-1 \to d=\text{"012"} \implies$ increment $i$.
        \item Else $\to d=\text{"210"} \implies$ increment $k$.
    \end{itemize}
    \item If $s=m-1$:
    \begin{itemize}
        \item If $i>0 \to d=\text{"120"} \implies$ increment $j$.
        \item Else $\to d=\text{"210"} \implies$ increment $k$.
    \end{itemize}
    \item Else ($0 < s < m-1$):
    \begin{itemize}
        \item If $i=m-1 \to d=\text{"201"} \implies$ increment $k$.
        \item Else $\to d=\text{"102"} \implies$ increment $j$.
    \end{itemize}
\end{itemize}

This construction is known to work for $c=0$ because it carefully manages the flow between the hyperplanes $s=0$ and $s=m-1$.
\begin{itemize}
    \item From $s=0$, the cycle always moves to $s=1$ (by incrementing $i$ or $k$).
    \item From $s=m-1$, the cycle always moves to $s=0$ (by incrementing $j$ or $k$).
    \item For intermediate $s$, the logic ensures a systematic sweep that covers the entire grid without repetition.
\end{itemize}
Empirical verification confirms this produces a single cycle of length $m^3$.

\section{Proof of Failure for $c=1$ and $c=2$}
We prove that the logic on page 2 \textbf{cannot} generate a Hamiltonian cycle for $c=1$ (and by symmetry, $c=2$).

\subsection{Analysis of Allowed Moves for $c=1$}
For cycle $c=1$, the action is determined by $d[1]$. We examine what moves are allowed in each case:

\begin{enumerate}
    \item \textbf{Case 1: $s = 0$.}
    The rules force $d \in \{\text{"012"}, \text{"210"}\}$.
    \begin{itemize}
        \item $\text{"012"}[1] = \text{'1'}$
        \item $\text{"210"}[1] = \text{'1'}$
    \end{itemize}
    \textbf{Result:} When $s=0$, the cycle \textit{must} increment $j$. It is \textbf{impossible} to increment $i$ or $k$.
    
    \item \textbf{Case 2: $s = m-1$.}
    The rules force $d \in \{\text{"120"}, \text{"210"}\}$.
    \begin{itemize}
        \item $\text{"120"}[1] = \text{'2'}$
        \item $\text{"210"}[1] = \text{'1'}$
    \end{itemize}
    \textbf{Result:} When $s=m-1$, the cycle \textit{must} increment $j$ or $k$. It is \textbf{impossible} to increment $i$.
    
    \item \textbf{Case 3: $0 < s < m-1$.}
    The rules force $d \in \{\text{"201"}, \text{"102"}\}$.
    \begin{itemize}
        \item $\text{"201"}[1] = \text{'0'}$
        \item $\text{"102"}[1] = \text{'0'}$
    \end{itemize}
    \textbf{Result:} When $0 < s < m-1$, the cycle \textit{must} increment $i$.
\end{enumerate}

\subsection{The Structural Contradiction}
A Hamiltonian cycle must visit all $m^3$ vertices. Consider the coordinate $i$. To visit all vertices, the cycle must increment $i$ frequently enough to reach $m-1$ and wrap around to $0$.

However, look at the constraints on changing $i$:
\begin{itemize}
    \item You \textbf{cannot} increment $i$ when $s=0$.
    \item You \textbf{cannot} increment $i$ when $s=m-1$.
    \item You \textbf{MUST} increment $i$ when $0 < s < m-1$.
\end{itemize}

This creates a "traffic jam" in the state space:
1.  The cycle can only change the $i$-coordinate when the sum $s$ is strictly between $0$ and $m-1$.
2.  As soon as the sum hits $s=0$ or $s=m-1$, the ability to change $i$ is revoked.
3.  Because the rules for $s=0$ and $s=m-1$ force moves in $j$ and $k$ only, the cycle gets trapped in a subset of the state space where $i$ is fixed.
4.  It cannot escape this subset to visit vertices with a different $i$-coordinate until $s$ returns to the "middle" zone.
5.  Due to the rigid nature of the rules, this "middle zone" is not accessible often enough, or the cycle enters a local loop within a specific $i$-slice before covering all $(j,k)$ pairs.

Consequently, the path length for $c=1$ is strictly less than $m^3$. It will visit a subset of vertices (often exactly $m^2$ or less, depending on $m$) and then repeat, failing to be Hamiltonian.

\section{Conclusion}
\begin{itemize}
    \item \textbf{Problem Status:} The problem of decomposing the graph into three Hamiltonian cycles \textbf{is solvable} for all odd $m$. This was confirmed by the original Python program mentioned in the paper.
    \item \textbf{Code Status:} The specific C code snippet on page 2 \textbf{is broken} for $c=1$ and $c=2$. It is an oversimplification that imposes contradictory constraints on the transition logic.
    \item \textbf{Implication:} The existence of a valid solution does not invalidate the claim that the problem is solved; it invalidates the \textit{provided implementation} of that solution. The true solution likely involves a more complex set of permutations or logic conditions not captured by the simplified snippet.
\end{itemize}

\end{document}

